% generated by roam.py handbook export -- не редактировать, править заметки
% 20260917205407-борелевская_оболочка.org (body af22c31c0e94)

\begin{definition}[{Борелевская оболочка}]\label[definition]{note:3819FAC4-B7DB-4181-A40B-8FBB5EC85DC1}
  Пусть $X$~--- множество, $\mathcal{E}$~--- система его подмножеств. \textbf{Борелевская оболочка} $\mathcal{B}(\mathcal{E})$~--- наименьшая $\sigma$-алгебра подмножеств $X$, содержащая $\mathcal{E}$:
  \[
  \mathcal{B}(\mathcal{E}) = \bigcap \left\{ \mathcal{A}\ \middle|\ \mathcal{A} \text{ -- } \sigma\text{-алгебра подмножеств } X,\ \mathcal{E} \subset \mathcal{A} \right\}
  \]

  Напоминание: система $\mathcal{A}$ подмножеств $X$~--- $\sigma$-алгебра, если

  \begin{enumerate}
    \item $X \in \mathcal{A}$;
    \item $A \in \mathcal{A} \Rightarrow X \setminus A \in \mathcal{A}$;
    \item $A_1, A_2, \dots \in \mathcal{A} \Rightarrow \bigcup\limits_{k=1}^{\infty} A_k \in \mathcal{A}$.
  \end{enumerate}
\end{definition}

\begin{remark}
  Главный пример~--- \emph{борелевская} $\sigma$-алгебра $\mathfrak{B}(\R^n)$: борелевская оболочка системы всех открытых множеств $\R^n$.
\end{remark}

\begin{remark}
  Та же схема «пересечь все объекты с нужным свойством, содержащие $M$» дает \hyperref[note:1AC657F9-2A3D-47D1-A71D-5B813D961260]{выпуклую оболочку}, линейную оболочку, замыкание множества. Она работает, когда (1) хотя бы один такой объект есть и (2) свойство сохраняется при любых пересечениях~--- как у $\sigma$-алгебр (ниже) и у \hyperref[note:08CF7113-4709-4C4F-A0E7-4A2ED595BD7F]{выпуклых множеств}.
\end{remark}

\begin{proof}
  Семейство в пересечении непусто: $2^X$ (все подмножества $X$)~--- $\sigma$-алгебра, содержащая $\mathcal{E}$.

  Пересечение $\mathcal{B}$ любого семейства $\sigma$-алгебр $\mathcal{A}_i$~--- снова $\sigma$-алгебра: каждая аксиома выполнена в каждой $\mathcal{A}_i$, поэтому и в пересечении. Например, если $A_1, A_2, \dots \in \mathcal{B}$, то все $A_k$ лежат в каждой $\mathcal{A}_i$, значит $\bigcup A_k \in \mathcal{A}_i$ для всех $i$, то есть $\bigcup A_k \in \mathcal{B}$.

  Минимальность: $\mathcal{E} \subset \mathcal{B}(\mathcal{E})$, и любая $\sigma$-алгебра $\mathcal{A} \supset \mathcal{E}$ участвует в пересечении, поэтому $\mathcal{B}(\mathcal{E}) \subset \mathcal{A}$.
\end{proof}
